\section{Approach}
\label{sec:approach}

\subsection{Overview}

Figure~\ref{fig:overview} shows the pipeline:
typed workflow generation $\rightarrow$ capability-aware lowering $\rightarrow$ parallel
execution over backend families $\rightarrow$ semantic normalization $\rightarrow$
compositional oracle evaluation $\rightarrow$ triage/classification $\rightarrow$ evidence
pipeline. \todo{Create Figure~\ref{fig:overview} (TikZ or PDF) showing the six stages and the
fresh/replay separation boundary.}

\subsection{Unified testing surface}
\label{sec:approach-surface}

\todo{Describe: (i) the typed workflow IR and small expression AST; (ii) the target registry
that declares each backend's family and capabilities; (iii) capability-aware lowering that
maps one workflow onto pandas/Polars APIs, Arrow compute, SQL, and query-engine execution;
(iv) the shared operation- and case-level semantics that keep generator, scheduler, oracle and
triage from drifting. Point to \texttt{src/datadiff/ccs_ir.py}, \texttt{dsl.py},
\texttt{targets.py}, \texttt{adapters/}, \texttt{backends/}.}

\subsection{Semantic normalization}
\label{sec:approach-normalizer}

\todo{Describe the canonical form: dtype promotion and canonical dtype names, null/NaN/negative
zero equivalence rules, index and column-order erasure where semantics allow, row ordering and
tie determinism, Arrow layout erasure (slice offset, chunking, dictionary encoding), numeric
tolerance policy, and error/timing classification. State explicitly which differences are
\emph{semantic} and which are \emph{presentational}. Point to \texttt{normalizer.py} and
\texttt{canonicalization.py}.}

\subsection{Compositional multi-endpoint oracle}
\label{sec:approach-oracle}

\todo{Describe: differential oracle across backend families; metamorphic relations on
semantics-preserving transformations (sort, filter/take, rebuild, pagination, string
normalization idempotence, coalesce idempotence, semi/anti join with duplicate keys); and
deterministic latest-version probes. Explain per-case fusion and confidence, the four verdicts,
and root-cause classification. \textbf{Recount probes and relations from the live registry
before publishing any number} (see \texttt{paper/CLAIMS\_RECONCILIATION.md} \S3).}

\subsection{Compositional semantic HyperContract (oracle core)}
\label{sec:approach-hypercontract}

\todo{Describe the phase-6 design: finite multi-endpoint contracts; forward semantic-property
inference and backward observability analysis; derivation / applicability / observation
certificates per observed obligation; staged planner vs authority exact comparator with
escalation; the declared target lattice (16 fresh families / 232 cells, 384 single-axis
contrast edges, 502 target-control pair obligations, 9 regression families / 144 cells); and
the coverage-debt scheduler with deterministic seed substreams. This is the paper's
software-engineering depth; keep it factual and separate unrun gates from completed ones.}

\subsection{Search and scheduling}

\todo{Describe typed generation profiles, mutation operators, the multi-scope adaptive
selection across case/profile/target/backend dimensions, the quality-diversity seed archive,
and lane scheduling. Frame as the mechanism behind RQ4; do not claim priority.}

\subsection{Evidence pipeline and counting contract}
\label{sec:approach-evidence}

\todo{Describe the five counting tiers (finding, recheck survivor, native candidate,
issue-ready root, strict confirmed family), family key = root cause + suspicious backends,
fresh/replay physical separation, reduction, issue bundling, and confirmation tracking. Point
to \texttt{triage.py}, \texttt{family\_novelty.py}, \texttt{reducer.py},
\texttt{issue\_bundle.py}, \texttt{final\_readiness.py}.}
